\documentclass[10pt,a4paper]{amsart}
\usepackage{amsmath,amssymb,amsthm,graphicx}
\usepackage{verbatim}
\newenvironment{reference}{\comment}{\endcomment}
%\newenvironment{reference}{}{}
\newenvironment{slogan}{\comment}{\endcomment}
\newenvironment{history}{\comment}{\endcomment}

% For commutative diagrams we use Xy-pic
\usepackage[all]{xy}

% We use 2cell for 2-commutative diagrams.
\xyoption{2cell}
% Theorem environments.
%
\theoremstyle{plain}
\newtheorem{theorem}[subsection]{Theorem}
\newtheorem{proposition}[subsection]{Proposition}
\newtheorem{lemma}[subsection]{Lemma}

\theoremstyle{definition}
\newtheorem{definition}[subsection]{Definition}
\newtheorem{example}[subsection]{Example}
\newtheorem{exercise}[subsection]{Exercise}
\newtheorem{situation}[subsection]{Situation}

\theoremstyle{remark}
\newtheorem{remark}[subsection]{Remark}
\newtheorem{remarks}[subsection]{Remarks}

\numberwithin{equation}{subsection}

% Macros
%
\def\lim{\mathop{\mathrm{lim}}\nolimits}
\def\colim{\mathop{\mathrm{colim}}\nolimits}
\def\Spec{\mathop{\mathrm{Spec}}}
\def\Hom{\mathop{\mathrm{Hom}}\nolimits}
\def\Ext{\mathop{\mathrm{Ext}}\nolimits}
\def\SheafHom{\mathop{\mathcal{H}\!\mathit{om}}\nolimits}
\def\SheafExt{\mathop{\mathcal{E}\!\mathit{xt}}\nolimits}
\def\Sch{\mathit{Sch}}
\def\Mor{\mathop{\mathrm{Mor}}\nolimits}
\def\Ob{\mathop{\mathrm{Ob}}\nolimits}
\def\Sh{\mathop{\mathit{Sh}}\nolimits}
\def\NL{\mathop{N\!L}\nolimits}
\def\CH{\mathop{\mathrm{CH}}\nolimits}
\def\proetale{{pro\text{-}\acute{e}tale}}
\def\etale{{\acute{e}tale}}
\def\QCoh{\mathit{QCoh}}
\def\Ker{\mathop{\mathrm{Ker}}}
\def\Im{\mathop{\mathrm{Im}}}
\def\Coker{\mathop{\mathrm{Coker}}}
\def\Coim{\mathop{\mathrm{Coim}}}

% Boxtimes
%
\DeclareMathSymbol{\boxtimes}{\mathbin}{AMSa}{"02}

%
% Macros for moduli stacks/spaces
%
\def\QCohstack{\mathcal{QC}\!\mathit{oh}}
\def\Cohstack{\mathcal{C}\!\mathit{oh}}
\def\Spacesstack{\mathcal{S}\!\mathit{paces}}
\def\Quotfunctor{\mathrm{Quot}}
\def\Hilbfunctor{\mathrm{Hilb}}
\def\Curvesstack{\mathcal{C}\!\mathit{urves}}
\def\Polarizedstack{\mathcal{P}\!\mathit{olarized}}
\def\Complexesstack{\mathcal{C}\!\mathit{omplexes}}
% \Pic is the operator that assigns to X its picard group, usage \Pic(X)
% \Picardstack_{X/B} denotes the Picard stack of X over B
% \Picardfunctor_{X/B} denotes the Picard functor of X over B
\def\Pic{\mathop{\mathrm{Pic}}\nolimits}
\def\Picardstack{\mathcal{P}\!\mathit{ic}}
\def\Picardfunctor{\mathrm{Pic}}
\def\Deformationcategory{\mathcal{D}\!\mathit{ef}}
\usepackage{cleveref}
\crefname{originalreference}{original reference}{original references}
\Crefname{originalreference}{Original reference}{Original references}
\makeatletter
\newlabel{definition-rig-etale-continuous-homomorphism}{{0GCV}{source}{}{sourcepointer-1}{}}
\newlabel{definition-rig-etale-continuous-homomorphism@cref}{{[originalreference][0][]0GCV}{[1][0][]source}{}{}{}}
\makeatother
\crefname{originalreference}{original reference}{original references}
\Crefname{originalreference}{Original reference}{Original references}
\makeatletter
\newlabel{definition-rig-etale-homomorphism}{{0GAY}{source}{}{sourcepointer-2}{}}
\newlabel{definition-rig-etale-homomorphism@cref}{{[originalreference][0][]0GAY}{[1][0][]source}{}{}{}}
\makeatother
\crefname{originalreference}{original reference}{original references}
\Crefname{originalreference}{Original reference}{Original references}
\makeatletter
\newlabel{definition-rig-smooth}{{0GCP}{source}{}{sourcepointer-3}{}}
\newlabel{definition-rig-smooth@cref}{{[originalreference][0][]0GCP}{[1][0][]source}{}{}{}}
\makeatother
\crefname{originallemma}{original lemma}{original lemmas}
\Crefname{originallemma}{Original lemma}{Original lemmas}
\makeatletter
\newlabel{formal-spaces-lemma-closed-immersion-into-countably-indexed}{{0AQI}{source}{}{sourcepointer-4}{}}
\newlabel{formal-spaces-lemma-closed-immersion-into-countably-indexed@cref}{{[originallemma][0][]0AQI}{[1][0][]source}{}{}{}}
\makeatother
\crefname{originalreference}{original reference}{original references}
\Crefname{originalreference}{Original reference}{Original references}
\makeatletter
\newlabel{lemma-equivalent-with-artin-smooth}{{0GAJ}{source}{}{sourcepointer-5}{}}
\newlabel{lemma-equivalent-with-artin-smooth@cref}{{[originalreference][0][]0GAJ}{[1][0][]source}{}{}{}}
\makeatother
\crefname{originalreference}{original reference}{original references}
\Crefname{originalreference}{Original reference}{Original references}
\makeatletter
\newlabel{lemma-rig-etale-morphisms}{{0GCZ}{source}{}{sourcepointer-6}{}}
\newlabel{lemma-rig-etale-morphisms@cref}{{[originalreference][0][]0GCZ}{[1][0][]source}{}{}{}}
\makeatother
\crefname{originalreference}{original reference}{original references}
\Crefname{originalreference}{Original reference}{Original references}
\makeatletter
\newlabel{lemma-rig-etale-rig-smooth-morphism}{{0GD0}{source}{}{sourcepointer-7}{}}
\newlabel{lemma-rig-etale-rig-smooth-morphism@cref}{{[originalreference][0][]0GD0}{[1][0][]source}{}{}{}}
\makeatother
\title{For closed immersions of locally Noetherian formal spaces, rig-smooth equals rig-etale}
\author{The Stacks Project Authors}
\date{stacks-project:restricted.tex:local-snapshot-20261012; GNU FDL 1.2 or later}
\begin{document}
\maketitle
\noindent\textit{Editorial note.} All delivered mathematical text is the author's own, copied from the pinned Stacks Project chapter file restricted.tex; no statement or proof step was rewritten, reordered or removed.\par\medskip
\noindent\textit{Editorial note.} The counted proof is the author's complete argument: the proof first identifies the ring maps corresponding to etale-local pieces of a closed immersion, then verifies the two implications through the description of NL\textasciicircum{}wedge as J/J\textasciicircum{}2 in degree -1, closing with the observation that a power of an ideal of definition is again an ideal of definition\par\medskip
\noindent\textit{Editorial note.} The label lemma-closed-immersion-rig-smooth names the counted statement, and the proof environment immediately following it in the source is delivered verbatim as its proof.\par\medskip
\noindent\textit{Editorial note.} The counted statement is carried by the author's own lemma environment in the chapter source.\par\medskip
\noindent\textit{Allowed standard prerequisites (editorial).}\begin{itemize}
\item closed immersions of locally countably indexed formal algebraic spaces correspond to surjective morphisms B -> A in WAdm\textasciicircum{}Noeth (Formal Spaces, Lemma formal-spaces-lemma-closed-immersion-into-countably-indexed)
\item rig-etale morphisms are rig-smooth (lemma-rig-etale-rig-smooth-morphism of the same chapter)
\item the definition of rig-etale morphisms of formal algebraic spaces via the ring maps of WAdm\textasciicircum{}Noeth (lemma-rig-etale-morphisms and definition-rig-etale-continuous-homomorphism of the same chapter)
\item the completed naive cotangent complex of a quotient B -> A = B/J is the complex with J/J\textasciicircum{}2 in degree -1 (definition-rig-etale-homomorphism of the same chapter)
\item rig-smooth ring maps and the criterion that annihilation of J/J\textasciicircum{}2 by a power of I follows from rig-smoothness (definition-rig-smooth and lemma-equivalent-with-artin-smooth of the same chapter)
\end{itemize}\medskip
\ifdefined\hypertarget\hypertarget{sourcepointer-1}{}\fi
\noindent\textit{Source reference.} Original reference 0GCV: sources/stacks-restricted/pinned-original-source.tex, source lines 4776--4776. Prerequisite pointer invoked at proof line(s) 5211; it names a standard fact of the same chapter that the counted argument uses, and every step of the counted proof built on it is retained verbatim in the delivered body.\par\medskip
\ifdefined\hypertarget\hypertarget{sourcepointer-2}{}\fi
\noindent\textit{Source reference.} Original reference 0GAY: sources/stacks-restricted/pinned-original-source.tex, source lines 1794--1794. Prerequisite pointer invoked at proof line(s) 5216; it names a standard fact of the same chapter that the counted argument uses, and every step of the counted proof built on it is retained verbatim in the delivered body.\par\medskip
\ifdefined\hypertarget\hypertarget{sourcepointer-3}{}\fi
\noindent\textit{Source reference.} Original reference 0GCP: sources/stacks-restricted/pinned-original-source.tex, source lines 4622--4622. Prerequisite pointer invoked at proof line(s) 5220; it names a standard fact of the same chapter that the counted argument uses, and every step of the counted proof built on it is retained verbatim in the delivered body.\par\medskip
\ifdefined\hypertarget\hypertarget{sourcepointer-4}{}\fi
\noindent\textit{Source reference.} Original lemma 0AQI: sources/stacks-chow/tags.txt, source lines 13511--13511. Prerequisite pointer invoked at proof line(s) 5200; it names a standard fact of a sibling chapter that the counted argument uses, and every step of the counted proof built on it is retained verbatim in the delivered body.\par\medskip
\ifdefined\hypertarget\hypertarget{sourcepointer-5}{}\fi
\noindent\textit{Source reference.} Original reference 0GAJ: sources/stacks-restricted/pinned-original-source.tex, source lines 843--843. Prerequisite pointer invoked at proof line(s) 5225; it names a standard fact of the same chapter that the counted argument uses, and every step of the counted proof built on it is retained verbatim in the delivered body.\par\medskip
\ifdefined\hypertarget\hypertarget{sourcepointer-6}{}\fi
\noindent\textit{Source reference.} Original reference 0GCZ: sources/stacks-restricted/pinned-original-source.tex, source lines 5047--5047. Prerequisite pointer invoked at proof line(s) 5208; it names a standard fact of the same chapter that the counted argument uses, and every step of the counted proof built on it is retained verbatim in the delivered body.\par\medskip
\ifdefined\hypertarget\hypertarget{sourcepointer-7}{}\fi
\noindent\textit{Source reference.} Original reference 0GD0: sources/stacks-restricted/pinned-original-source.tex, source lines 5114--5114. Prerequisite pointer invoked at proof line(s) 5204; it names a standard fact of the same chapter that the counted argument uses, and every step of the counted proof built on it is retained verbatim in the delivered body.\par\medskip
\begin{lemma}
\label{lemma-closed-immersion-rig-smooth}
Let $S$ be a scheme. Let $f : X \to Y$ be a morphism of formal algebraic
spaces. Assume that $X$ and $Y$ are locally Noetherian and $f$ is a
closed immersion. The following are equivalent
\begin{enumerate}
\item $f$ is rig-smooth,
\item $f$ is rig-\'etale,
\item for every affine formal algebraic space $V$ and every morphism
$V \to Y$ which is representable by algebraic spaces and \'etale
the morphism $X \times_Y V \to V$ corresponds to a surjective morphism
$B \to A$ in $\textit{WAdm}^{Noeth}$ whose kernel $J$ has the following
property: $I(J/J^2) = 0$ for some ideal of definition $I$ of $B$.
\end{enumerate}
\end{lemma}

\begin{proof}
Let us observe that given $V$ and $V \to Y$ as in (2) without any
further assumption on $f$ we see that the morphism $X \times_Y V \to V$
corresponds to a surjective morphism $B \to A$ in $\textit{WAdm}^{Noeth}$
by Formal Spaces, Lemma
\ref{formal-spaces-lemma-closed-immersion-into-countably-indexed}.

\medskip\noindent
We have (2) $\Rightarrow$ (1) by
Lemma \ref{lemma-rig-etale-rig-smooth-morphism}.

\medskip\noindent
Proof of (3) $\Rightarrow$ (2). Assume (3). By
Lemma \ref{lemma-rig-etale-morphisms}
it suffices to show that the ring maps
$B \to A$ occurring in (3) are rig-\'etale in the
sense of Definition \ref{definition-rig-etale-continuous-homomorphism}.
Let $I$ be as in (3). The naive cotangent complex
$\NL_{A/B}^\wedge$ of $A$ over $(B, I)$ is the complex of $A$-modules
given by putting $J/J^2$ in degree $-1$. Hence $A$ is
rig-\'etale over $(B, I)$ by
Definition \ref{definition-rig-etale-homomorphism}.

\medskip\noindent
Assume (1) and let $V$ and $B \to A$ be as in (3).
By Definition \ref{definition-rig-smooth} we see that
$B \to A$ is rig-smooth. Choose any ideal of definition $I \subset B$.
Then $A$ is rig-smooth over $(B, I)$.
As above the complex $\NL_{A/B}^\wedge$ is 
given by putting $J/J^2$ in degree $-1$.
Hence by Lemma \ref{lemma-equivalent-with-artin-smooth}
we see that $J/J^2$ is annihilated by
a power $I^n$ for some $n \geq 1$. Since $B$ is adic, we see
that $I^n$ is an ideal of definition of $B$ and the
proof is complete.
\end{proof}

\end{document}
